\section{Dashboard Creation}

Dashboards transform raw telemetry into visual insights that help analysts monitor activity, validate detections, and identify anomalies at scale. For a CTI analyst or threat hunter, dashboards serve two purposes: they provide situational awareness, and they act as investigative starting points. A well‑designed dashboard highlights deviations from normal behavior, exposes suspicious trends, and allows analysts to pivot quickly into deeper searches. This section introduces the core principles of dashboard design and provides practical SPL and KQL examples that can be used to build effective visualizations in Splunk and Microsoft Sentinel.

A good dashboard begins with a clear objective. Instead of trying to visualize everything, focus on the behaviors that matter most: process execution, authentication patterns, network activity, and high‑risk events. Each panel should answer a specific question, such as “Which processes are most active?”, “Which accounts are failing logons?”, or “Which hosts are making rare outbound connections?” By structuring dashboards around investigative questions, you ensure that each visualization provides actionable insight rather than noise.

\subsection*{Process Activity Panel}

Process creation is one of the most valuable signals for both hunting and detection engineering. A dashboard panel that highlights the most common processes helps establish a baseline and makes anomalies easier to spot.

\textbf{SPL Example}
\begin{lstlisting}
index=sysmon EventCode=1
| stats count by Image
| sort - count
\end{lstlisting}

\textbf{KQL Equivalent}
\begin{lstlisting}
Sysmon
| where EventID == 1
| summarize Count=count() by Image
| order by Count desc
\end{lstlisting}

This panel helps analysts quickly identify unusual binaries or unexpected spikes in activity.

\subsection*{Authentication Trends Panel}

Authentication logs reveal credential misuse, brute‑force attempts, and lateral movement. A dashboard panel that tracks failed logons over time provides early warning of suspicious behavior.

\textbf{SPL Example}
\begin{lstlisting}
index=wineventlog EventCode=4625
| timechart count by AccountName
\end{lstlisting}

\textbf{KQL Equivalent}
\begin{lstlisting}
SecurityEvent
| where EventID == 4625
| summarize Count=count() by bin(TimeGenerated, 1h), Account
| render timechart
\end{lstlisting}

This visualization helps identify targeted accounts or sudden increases in authentication failures.

\subsection*{Network Activity Panel}

Network telemetry is essential for detecting command‑and‑control activity, lateral movement, and exfiltration. A dashboard panel that highlights rare outbound connections can reveal early signs of compromise.

\textbf{SPL Example}
\begin{lstlisting}
index=sysmon EventCode=3
| stats count by DestinationIp, DestinationPort
| where count < 5
\end{lstlisting}

\textbf{KQL Equivalent}
\begin{lstlisting}
Sysmon
| where EventID == 3
| summarize Count=count() by DestinationIp, DestinationPort
| where Count < 5
\end{lstlisting}

This panel surfaces low‑frequency connections that may indicate malicious activity.

\subsection*{High‑Risk Event Panel}

Some events are inherently high‑risk and deserve dedicated visibility. Examples include PowerShell execution, scheduled task creation, and privilege escalation attempts. A dashboard panel that highlights these events helps analysts quickly identify suspicious behavior.

\textbf{SPL Example}
\begin{lstlisting}
index=sysmon (EventCode=1 Image="*powershell.exe") OR EventCode=13 OR EventCode=7
| stats count by EventCode, Image, User
\end{lstlisting}

\textbf{KQL Equivalent}
\begin{lstlisting}
Sysmon
| where EventID in (1, 7, 13)
| summarize Count=count() by EventID, Image, User
\end{lstlisting}

This panel consolidates multiple high‑risk signals into a single view.

\subsection*{Design Principles}

Effective dashboards follow a few simple principles. Panels should be arranged from broad to specific, allowing analysts to move from general trends to detailed insights. Visualizations should be easy to interpret at a glance, with clear labels and minimal clutter. Most importantly, dashboards should be actionable: every panel should support a pivot into deeper investigation. Over time, dashboards should be refined based on analyst feedback and evolving threat behavior.

The examples in this section provide the foundation for building dashboards that support real‑world CTI, hunting, and detection workflows. Later sections will expand on these concepts with full dashboard layouts and advanced visualizations tailored to specific adversary techniques.
